\uuid{6714}
\titre{Notion d'équation différentielle}
\auteur{bodin ; croizat ; sacre}
\organisation{exo7}
\licence{CC BY-NC-SA 4.0}
\datecreate{2022-03-08}
\niveau{L1}
\module{Analyse}
\chapitre{Equation différentielle}
\sousChapitre{Résolution d'équation différentielle du premier ordre}
\difficulte{2}
\theme{Notion d'équation différentielle}
\exo7qcmid{hilisit.equations-differentielles.018}

\contenu{
\begin{question}
\qclassification{225.01}
Soit l'équation différentielle $xy'-3y=0$. Quelles sont les affirmations vraies ?
\begin{answers}
  \bad{$x^3+1$ est une solution.}
  \good{$x^3$ est une solution.}
  \bad{$ e^{3x}$ est une solution.}
  \good{La fonction nulle est la seule solution constante.}
\end{answers}
\begin{explanations}
Pour une solution constante $y=C$, $y'=0$ donc $3y=0$ donc $y$ est la fonction nulle (et réciproquement, la fonction nulle est bien solution). Pour $y=x^3$, $xy'-3y=x \cdot  3x^2-3x^3=0$ donc $x^3$ est solution. Pour $y=x^3+1$, $xy'-3y=x \cdot 3x^2-3x^3-3=-3$ donc $x^3+1$ n'est pas solution. Pour $y= e^{3x}$, $xy'-3y=x \cdot 3 e^{3x}-3 e^{3x}=3(x-1) e^{3x}$, ce qui n'est pas la fonction nulle, donc $y= e^{3x}$ n'est pas solution.
\end{explanations}
\end{question}
}
